\documentclass[10pt]{article}

\usepackage[utf8]{inputenc}

\usepackage[T1]{fontenc}

\usepackage{amsmath}

\usepackage{amsfonts}

\usepackage{amssymb}

\usepackage[version=4]{mhchem}

\usepackage{stmaryrd}

\usepackage{bbold}



\begin{document}

имеет вид:



\[

u(x)=\int_{\mathbb{R}_{0}^{n}} f(y) P \mathbb{R}_{+}^{n}(x, y) d \mathbb{R}_{0}^{n}(y),

\]



где



\[

\mathbb{R}_{0}^{n}=\left\{y=\left(y_{1}, \ldots, y_{n}\right) \in \mathbb{R}^{n}: y_{n}=0\right\},

\]



$d \mathbb{R}_{0}^{n}$ - элемент площади $\mathbb{R}_{0}^{n}, f(y)$ - ограниченная непрерывная функция на $\mathbb{R}_{0}^{n}$,



\[

P \mathbb{R}_{+}^{n}(x, y)=\frac{2 \Gamma\left(\frac{n}{2}+1\right)}{n \pi^{n / 2}} \frac{x_{n}}{|x-y|^{n}}

\]



\begin{itemize}

  \item ядро Пуассона для полупространства. Формулы (1) и (3) суть частные случаи формулы Грина:

\end{itemize}



\[

u(x)=\int_{\Gamma} f(y) \frac{\partial G(x, y)}{\partial n_{y}} d \Gamma(y)

\]



дающей решение задачи Дирихле для областей $D \subset \mathbb{R}^{n}$ с гладкой границей Г при помощи производной $d G(x, y) / d n_{y}$ функции Грина $G(x, y)$ по направлению внутренней нормали к Г в точке $y \in \Gamma$. Иногда формулу (4) также называют интегралом Пуассона.



Основные свойства П. и.: 1) $u(x)$ есть гармоническая функция координат точки $x$; 2) П. и. дает решение задачи Дирихле с граничными данными $f(y)$ в классе (ограниченных) гармонич. функций, то есть функция $u(x)$, продолженная на границу области значениями $f(y)$, непрерывна в замкнутой области. На этих свойствах основаны применения П. и. в классич. математич. физике.



Большую роль в теории аналитич. функций многих комплексных переменных и в ее применениях к квантовой теории поля играют различные модификации П. и. Напр., ядро Пуассона для поликруга



\[

U^{n}=\left\{z=\left(z_{1}, \ldots, z_{n}\right) \in \mathbb{C}^{n}:\left|z_{j}\right|<1, j=1, \ldots, n\right\}

\]



комплексного пространства $\mathbb{C}^{n}$ получается при перемножении ядер (2):



\[

P U^{n}(z, \xi)=\prod_{j=1}^{n} P B_{2}\left(z_{j}, \xi_{j}\right)

\]



Соответствуюший П.и.



\[

u(z)=\int_{T^{n}} f(\xi) P U^{n}(z, \xi) d T^{n}(\xi)

\]



по остову поликруга



\[

T^{n}=\left\{\zeta_{1}\left(\zeta_{1}, \ldots, \zeta_{n}\right) \in \mathbb{C}^{n}:\left|\zeta_{j}\right|=1, j=1, \ldots, n\right\}

\]



дает кратногармонич. функцию $u(z), z \in U^{n}$, принимающую на остове $T^{n}$ непрерывные значения $f(\zeta)$.



В квантовой теории поля применяются П. и. для трубчатых областей $T^{C}$ комплексного пространства $\mathbb{C}^{n}$ над выпуклым открытым острым конусом $C$ в пространстве $\mathbb{R}^{n}$ (с вершиной в начале координат) вида:



\[

\begin{aligned}

& T^{C}=\mathbb{R}^{n}+i \mathbb{C}=\left\{z=x+i y \in \mathbb{C}^{n} ;\right. \\

& \left.\quad x=\left(x_{1}, \ldots, x_{n}\right) \in \mathbb{R}^{n} ; y=\left(y_{1}, \ldots, y_{n}\right) \in \mathbb{C}\right\} .

\end{aligned}

\]



П. и. цля полуплоскости вида (3) при $n=2$ есть частный случай таких П. и. для трубчатых областей.



Лит.: [1] Т и х о н о в А. Н., С а м а р с к и й А. А., Уравнения математической физики, 5 изд., М., $1977 . \quad$ Е.Д. Соломенцев.



ПУАССОНА РАСПРЕДЕЛЕНИЕ - распределение вероятностей случайной величины $N$, принимающей целые неотрицательные значения $n=0,1,2, \ldots$ с вероятностями



\[

P_{n}=\mathrm{P}\{N=n\}=\frac{\lambda^{n}}{n !} e^{-\lambda}

\]



где $\lambda$ - параметр распределения. Характеристич. функция П. р. имеет вид:



\[

f(t)=\mathrm{E} e^{i N t}=\exp \left[\lambda\left(e^{i t}-1\right)\right]

\]



и поэтому семиинварианты любого порядка равны $\lambda$.



Сумма независимых случайных величин $N_{1}, N_{2}, \ldots, N_{s}$, имеющих П. р. с параметрами $\lambda_{1}, \lambda_{2} \ldots, \lambda_{s}$, подчиняется П. р. с параметром $\left(\lambda_{1}+\lambda_{2}+\ldots+\lambda_{s}\right)$ и, обратно, если сумма $\left(N_{1}+N_{2}\right)$ двух независимых случайных величин $N_{1}$ и $N_{2}$ имеет П. р., то каждая случайная величина $N_{1}$ и $N_{2}$ подчинена П. р.



П. р. было впервые получено С. Пуассоном (S. Poisson, 1837) как асимптотич. распределение вероятностей числа «успехов» $N$ в схеме Бернулли при большом числе испытаний $M$ и малой вероятности успеха $p$ :



\[

\mathrm{P}\{N=n\}=C_{M}^{n} p^{n}(1-p)^{M-n} \rightarrow \frac{\lambda^{n}}{n !} e^{-\lambda},

\]



при $M \rightarrow \infty, p \rightarrow 0, M p \rightarrow \lambda$. Это обстоятельство обуславливает то, что П. р. является хорошим приближением при моделировании многих явлений, связанных с появлением случайного числа редких слабозависимых событий в фиксированном промежутке времени и в ограниченной области пространства. Наряду с обыкновенным П. р. в математич. физике используется и так наз. обобще нное, или сложное, П. р. для вероятностей значений суммы $\left(\xi_{1}+\xi_{2}+\ldots+\xi_{N}\right)$ случайного числа $N$ одинаково распределенных случайных величин $\xi_{1}, \xi_{2}, \ldots$. При этом $N ; \xi_{1}, \xi_{2}, \ldots-$ взаимно независимы и $N$ распределено согласно П. р. с параметром $\lambda$. Характеристич. функция $f(t)$ сложного П. р. имеет вид:



\[

f(t)=\exp [\lambda(\psi(t)-1)]

\]



где $\psi(t)$ - характеристич. функция $\xi_{s}$



Если распределение вероятностей целой неотрицательной случайной величины $N$ определяется случайной величиной $\xi$ так, что условное распределение вероятностей $\mathrm{P}\{N=n \mid \xi=x\}$ является П. р. с параметром $\xi=x$, то оно называется с о с тавны м П. р. и определяется формулой



\[

P_{n}=\mathrm{P}\{N=n\}=\frac{1}{n !} \int_{0}^{\infty} d x p(x) x^{n} e^{-x},

\]



где $p(x)$ - плотность распределения величины $\xi$. Характеристй. функция составного П. р. имеет вид:



\[

f(t)=\mathrm{E} e^{i t N}=\varphi\left(i\left(1-e^{i t}\right)\right),

\]



где $\varphi(t)=\mathrm{E} e^{i t \xi}$. Составное П. р. находит применение в задачах квантовой оптики.



Лит.: [1] Ш и р я е в А. Н., Вероятность, М., 1980; [2] Бо ров к о в А. А., Теория вероятностей, 2 изд., М., 1986; [3] Фел л е р В., Введение в теорию вероятностей и ее приложения, пер. с англ., т. 1 , 2 изд., М., 1984; [4] Колч и н В. Ф., Се вастьян о в Б. А., Чист я к о в В.П., Случайные размещения, М., 1980. Ю. П. Вирченко. ПУАССОНА СКОБКИ - важное понятие аналитической механики, введенное С. Пуассоном (S. Poisson, 1809) и получившее дальнейшее развитие в гамильтоновой механике (см. Гамильтонов формализм). П. с. могут быть обобщены на случай квантовой механики, а также классич. и квантовой теории поля. С кобко й П у а с о на двух динамич. величин $f$ и $g$ нек-рой гамильтоновой системы называют выражение



\[

\{f, g\}=\sum_{k=1}^{n}\left(\frac{\partial f}{\partial p_{k}} \frac{\partial g}{\partial q_{k}}-\frac{\partial f}{\partial q_{k}} \frac{\partial g}{\partial p_{k}}\right)

\]



где $f(q, p, t)$ и $g(q, p, t)$ - нек-рые функции так наз. гамильтоновых (канонических) переменных $q_{1}, \ldots, p_{n}(n-$ число степеней свободы системы). Встречается определение П. с. $\{f, g\}$, отличающееся от определения (1) множителем $(-1)$. Для обозначения П. с. могут использоваться также круглые $(f, g)$ или квадратные $[f, g]$ скобки. Иногда термин употребляется в единственном числе - «скобка Пуассона». Из определения (1) следуют свойства П. с.:



\[

\begin{gathered}

\{g, f\}=-\{f, g\} ; \\

\{\alpha f+\beta g, h\}=\alpha\{f, h\}+\beta\{g, h\},

\end{gathered}

\]



где $\alpha, \beta$ - нек-рые константы;



ПУАССОНА 473





\end{document}